\section{Драйверы в операционной системе Windows}
\label{sec:Chapter3} \index{Chapter3}

Третья глава рассматривает архитектуру драйверов Windows и способы доступа
пользовательских и режимных ядра компонентов к разделяемому кольцевому буферу,
инициализированному прошивкой UEFI. Драйвер Windows завершает трёхкомпонентную
схему канала диагностики «гость --- Stratovirt», описанную во введении.

\subsection{Роль драйвера режима ядра}

В гостевой Windows диагностические сообщения генерируют:
\begin{itemize}
    \item пользовательские приложения и службы;
    \item другие драйверы режима ядра;
    \item компоненты, уже использующие ETW или \texttt{DbgPrint}, если их
    перенаправляют на общий API записи.
\end{itemize}

Драйвер режима ядра выступает единой точкой доступа к кольцевому буферу в
физической памяти гостя: он отображает область, выполняет проверки целостности
и уведомляет виртуальное устройство Stratovirt через MMIO-регистр \texttt{KICK}
(глава~\ref{sec:Chapter1}). Это исключает прямой доступ user-mode к сырым
физическим адресам и централизует политику безопасности.

\subsection{Архитектура ввода-вывода Windows}

Модель Windows разделяет выполнение на пользовательский режим и режим ядра.
Запросы к устройствам оформляются как IRP (I/O Request Packet), которые
проходят стек драйверов, начиная с I/O Manager~\cite{windows_wdk}.

Для виртуального устройства без «железного» IRQ типична схема:
\begin{enumerate}
    \item PnP Manager обнаруживает устройство по ACPI;
    \item функциональный драйвер создаёт \texttt{DEVICE\_OBJECT};
    \item пользовательский код открывает символьную ссылку на устройство
    (\texttt{\textbackslash{}.\textbackslash{}LogDev});
    и вызывает \texttt{DeviceIoControl};
    \item драйвер обрабатывает IOCTL и обращается к разделяемой памяти.
\end{enumerate}

\subsection{Модель разрабатываемого драйвера}

В реализации используется классическая модель WDM (Windows Driver Model):
точка входа \texttt{DriverEntry}, регистрация \texttt{DriverUnload},
создание устройства через \texttt{IoCreateDevice}, установка флагов
\texttt{DO\_BUFFERED\_IO} или \texttt{DO\_DIRECT\_IO} в зависимости от
способа передачи IOCTL-буферов.

Упрощённая сигнатура точки входа:

\begin{lstlisting}
NTSTATUS DriverEntry(
    PDRIVER_OBJECT DriverObject,
    PUNICODE_STRING RegistryPath
);
\end{lstlisting}

В \texttt{DriverEntry} регистрируются:
\begin{itemize}
    \item \texttt{MajorFunction[IRP\_MJ\_CREATE]} --- открытие устройства;
    \item \texttt{MajorFunction[IRP\_MJ\_CLOSE]} --- закрытие;
    \item \texttt{MajorFunction[IRP\_MJ\_DEVICE\_CONTROL]} --- обработка IOCTL;
    \item \texttt{DriverUnload} --- отмена отображения памяти и удаление устройства.
\end{itemize}

Альтернатива --- KMDF (создание очереди запросов \texttt{WdfIoQueue}), однако
для устройства с фиксированным MMIO и ACPI-ресурсом WDM остаётся прозрачным
и достаточным.

\subsection{Обнаружение буфера через ACPI}
\label{sec:acpi-logdev}

Критично, что физический адрес и размер кольцевого буфера не зашиты в драйвер
константами: их задаёт прошивка при выделении runtime-области, а Stratovirt
публикует MMIO-интерфейс в ACPI.

Эмулятор добавляет в таблицы ACPI устройство~\cite{stratovirt_repo}:
\begin{itemize}
    \item имя ACPI: \texttt{LOGD};
    \item \texttt{\_HID}: строка \texttt{STVRT\_LOG};
    \item \texttt{\_CRS}: дескриптор \texttt{Memory32Fixed} для MMIO-региона
    LogDev (база и длина совпадают с картой standard VM, например
    \texttt{0x0903\_0000}, размер \texttt{0x30}).
\end{itemize}

Драйвер Windows на этапе \texttt{IRP\_MN\_START\_DEVICE} получает список
ресурсов \texttt{CM\_RESOURCE\_LIST}. Из него извлекается диапазон MMIO для
регистров \texttt{BUFPTR}, \texttt{KICK} и др. Адрес кольцевого буфера в RAM
гостя считывается записью в регистр \texttt{BUFPTR} на этапе инициализации
(после того как UEFI записал туда физический адрес) либо передаётся через
дополнительный ACPI-метод / таблицу конфигурации платформы, если она
добавлена в образ прошивки.

Далее драйвер отображает физическую страницу (или непрерывный диапазон)
в системное виртуальное адресное пространство ядра функциями
\texttt{MmMapIoSpace} / \texttt{MmGetPhysicalAddress} (в зависимости от
типа памяти) и сохраняет указатель на структуру заголовка буфера, совместимую
с прошивочным layout (раздел~\ref{sec:Chapter4}).

Использование \texttt{\_HID = STVRT\_LOG} позволяет однозначно отличить
виртуальную платформу Stratovirt от физической машины, где такого устройства
нет, и не загружать драйвер на реальном железе.

\subsection{Интерфейс IOCTL для пользовательского режима}

Для приложений и служб user-mode предусмотрен IOCTL, принимающий готовое
сообщение (NUL-терминированная строка или структура с полем длины).
Обработчик в \texttt{IRP\_MJ\_DEVICE\_CONTROL}:
\begin{enumerate}
    \item проверяет входной буфер (\texttt{IoGetCurrentIrpStackLocation});
    \item валидирует размер (не превышает максимального сообщения кольца);
    \item вызывает общую функцию записи в lock-free буфер (та же, что в UEFI);
    \item при успехе записывает в MMIO \texttt{KICK} значение \texttt{0xDEADBABA}
    для пробуждения потока Stratovirt.
\end{enumerate}

Коды IOCTL задаются через \texttt{CTL\_CODE} с уникальным \texttt{DeviceType}
и функциями \texttt{METHOD\_BUFFERED} (копирование буфера ядром) либо
\texttt{METHOD\_NEITHER} при жёстком контроле Probe в драйвере~\cite{windows_wdk}.

\subsection{Интерфейс для драйверов режима ядра}

Второй IOCTL (или экспорт через \texttt{IoCreateDevice} + интерфейс
\texttt{GUID}) возвращает указатель на функцию записи в буфер:
\begin{lstlisting}
typedef NTSTATUS (*LOGDEV_WRITE_FN)(
    PVOID Message,
    ULONG MessageSize
);
\end{lstlisting}

Другой драйвер, получив указатель при инициализации, может логировать без
прохождения user-kernel transition и без повторного отображения памяти.
Это снижает задержку при высокой частоте событий от сетевого или дискового
стека.

Проверка подписи вызывающего модуля (опционально) выполняется через
сравнение уровня IRQL и принадлежности к загруженным драйверам; в учебной
реализации достаточно ограничить IOCTL привилегированным открытием устройства.

\subsection{Взаимодействие с Stratovirt после записи}

Последовательность для асинхронного режима совпадает с UEFI:
\begin{enumerate}
    \item \texttt{reserve} --- резервирование места в кольце (атомарные
    операции над \texttt{ShWPos}, см.~раздел~\ref{sec:Chapter4});
    \item копирование тела сообщения;
    \item \texttt{commit} --- обновление \texttt{WritePos} при нулевом счётчике
    активных писателей;
    \item запись в MMIO \texttt{KICK} или явный \texttt{FLUSH}.
\end{enumerate}

Поток \texttt{iothread} на хосте считывает сообщения и печатает их в общий
лог Stratovirt с префиксом \texttt{GUEST(ptr mode)}, что позволяет объединить
их с логами эмулятора и с сообщениями UEFI, как указано во введении.

\subsection{Сравнение путей доступа}

\begin{table}[h]
    \centering
    \caption{Способы записи диагностики в гостевой Windows}
    \label{tab:win-paths}
    \begin{tabular}{|l|p{4cm}|p{5cm}|}
        \hline
        Клиент & API & Особенности \\
        \hline
        Приложение / служба & \texttt{CreateFile} + IOCTL &
        Проверка буфера в ядре, стандартная модель безопасности \\
        Другой драйвер & указатель \texttt{LOGDEV\_WRITE\_FN} &
        Минимальная задержка, без syscall \\
        Прошивка (до ОС) & DebugLib / MMIO &
        Описано в главе~\ref{sec:Chapter2} \\
        \hline
    \end{tabular}
\end{table}

\subsection{Выводы по главе}

\begin{enumerate}
    \item Драйвер Windows связывает user-mode и kernel-mode клиентов с тем же
    кольцевым буфером, что инициализирует UEFI.
    \item ACPI \texttt{STVRT\_LOG} / \texttt{LOGD} обеспечивает обнаружение
    MMIO-устройства и согласованность с Stratovirt.
    \item Два IOCTL (запись сообщения и получение указателя на функцию записи)
    покрывают основные сценарии диагностики в госте.
    \item Детали lock-free алгоритма и структуры буфера приведены
    в разделе~\ref{sec:Chapter4}.
\end{enumerate}
